%%%%%%%%%%%%%%%%%%%%%%%%%%%%%%%%%%%%%%%%%%%%%%%%%
% Chapter: Reserved Keys
%%%%%%%%%%%%%%%%%%%%%%%%%%%%%%%%%%%%%%%%%%%%%%%%%
\chapter{Reserved Keys}
\label{chap:api_rsvd_keys}


Reserved keys are keys whose string representation begin with a prefix of
``\code{"pmix"}''. By definition, reserved keys are provided by the host
environment and the \ac{PMIx} server and required to be available at client
start of execution. \ac{PMIx} clients and tools are therefore prohibited from
posting reserved keys.

Host environments and \ac{PMIx} implementations may choose to define their own
custom-prefixed keys - if defined, these are to be treated as reserved. Users
are advised to check their implementation and host environment documentation
for a list of any custom prefixes they must avoid.

The \ac{PMIx} standard defines a relatively small set of APIs and the caller may customize the behavior of the API by passing one or more attributes to that API.
Additionally, attributes may be keys passed to \refapi{PMIx_Get} calls to access the specified values from the system.

Each attribute is represented by a \var{key} string, and a type for the associated \var{value}.
This section defines a set of \textbf{reserved} keys which are prefixed with \code{pmix.} to designate them as PMIx standard reserved keys. All definitions were introduced in version 1 of the standard unless otherwise marked.

Applications or associated libraries (e.g., \ac{MPI}) may choose to define additional attributes.
The attributes defined in this section are of the system and job as opposed to the attributes that the application (or associated libraries) might choose to expose.
Due to this extensibility the \refapi{PMIx_Get} API will return \refconst{PMIX_ERR_NOT_FOUND} if the provided \var{key} cannot be found.

Attributes added in this version of the standard are shown in \textit{\textbf{\color{magenta}magenta}} to distinguish them from those defined in prior versions, which are shown in \textit{\textbf{black}}. Deprecated attributes are shown in \textit{\textbf{\color{green!80!black}green}} and may be removed in a future version of the standard.

\declareAttribute{PMIX_ATTR_UNDEF}{NULL}{NULL}{
Constant representing an undefined attribute.
}

%
\declareAttributeNEW{PMIX_SERVER_INFO_ARRAY}{"pmix.srv.arr"}{pmix_data_array_t}{
Array of data on a given server, starting with its \refattr{PMIX_NSPACE}
}

%%%%%%%%%%%%%%%%%%%%%%%%%%%%%%
\section{Instant On}

The ``instant on'' concept is captured by the following definition:

\begin{addmargin}[2em]{2em}
\emph{All information required for setup and communication (including the address vector of endpoints for every process) is available to each process at start of execution}
\end{addmargin}

%Stage 4. During their initialization, each process connects to its local PMIx server by calling the PMIx_Init function. The PMIx server notifies the host RM of the connection, and returns to the process all information (or a pointer to shared memory where the data resides) provided by the local RM daemon during registration.
%At this point, the launch sequence is complete and application processes are free to communicate. While PMIx provides interfaces by which application processes can globally exchange data amongst themselves (e.g., BCX), these interfaces are primarily intended for use by legacy systems and are discouraged for future systems. Focus instead is being placed on enabling exchange-free operations.

%Note that some operations may always require inter-process exchange of information. For example, shared memory com- munication typically requires that one process create a backing file and communicate the name and location of that file to its local peers. In cases such as this where only a local exchange is involved, libraries can use the PMIX_LOCAL scope when posting the information to restrict the information to the set of processes executing on the same node, thus avoiding any cross-node exchange. As of this writing, the PMIx community has not identified any non-local scope information that cannot be provided by the host RM.

%Fully removing the BCX and initialization time barrier from the PMIx launch sequence requires modifications to the fabric manager and Network Interface Card (NIC) libraries, as well as potential NIC hardware changes for optimized performance. The PMIx community is working with network vendors to field this support in their product roadmaps, including support for:
%• host address resolution. Historically, each application process would independently discover the locally available NICs and determine their addressing information (e.g., TCP addresses or GID/LID pairs) and “broadcast” that information to its peers. While it is sometimes possible to resolve NIC addresses given a hostname, failure to provide addressing informa- tion up front can lead to non-scalable hostname resolution demands on the SMS resolving agent, thereby necessitating the all-to-all broadcast. Removing this exchange is accomplished by having the fabric manager provide any required ad- dressing information for all nodes in the allocation so the information can be included in the initial allocation setup message in Stage 3.
%• pre-launch assignment of rendezvous endpoints (e.g., a socket or queue pair index) to each process for initiating peer- to-peer communications. This does not preclude the dynamic assignment of additional endpoints (e.g., multiple queue

%pairs for enhanced bandwidth), but instead provides a means by which the two peers can initially contact each other to further establish their communication channel. Note that there is no restriction preventing the two peers from simply using the provided endpoint for regular communications if they so choose.
%unexpected messages (i.e. messages received prior to process registration with the receiving NIC). Since a process is given all required information to communicate with any peer upon startup, it is possible that a process could immediately at- tempt to send a message to a remote peer that has not yet been started. Network libraries currently can lose messages and/or return an error on communications addressed to a process that is not yet known to the receiving library. Thus, current parallel programming libraries typically include an out-of-band barrier operation just prior to enabling commu- nications to ensure that all processes have been instantiated and have registered with their local NICs. This requirement can be eliminated if sufficient retry logic and buffer space is included in the network library and/or NIC hardware. Avoiding aggressive retries may be desirable — e.g., the receiving NIC could return a specific error code indicating that the intended recipient is not known to the NIC, thus instructing the sending NIC to introduce some slight (possibly ran- domized) delay prior to retrying — to mitigate unnecessary network congestion during startup. Current measurements, however, indicate that the time difference between startup of processes on even an exascale machine would be in the tenth-of-a-second range, and non-benchmark applications typically do not immediately send to their peers — therefore, the value of taking extraordinary measures remains to be seen.
%Achieving all three of these steps will take some time and effort. In the interim, and in situations where fabrics cannot provide the required endpoint information prior to application launch, PMIx offers two benefits that can be exploited to improve launch performance. First, the amount of data included in the BCX can, in many cases, be substantially reduced by removing data provided by the local host RM to each process. For example, programming libraries often have their processes post the hostname of their location and their assigned binding. While the information from each process may involve a relatively small number of bytes, the aggregate total flowing through the collective of a large scale application can become quite large. Reviewing the code to avoid inclusion of data in the BCX that is already available via PMIx has been used in such cases to reduce the BCX collective message and subsequent launch time.

%%%%%%%%%%%%%%%%%%%%%%%%%%%%%%
\section{Data realms}
\label{api:struct:attributes:retrieval}

\ac{PMIx} information spans a wide range of topics and sources. In some cases,
there are multiple overlapping sources for the same type of data - e.g., the
session, job, and application can each provide information on the number of
nodes involved in the respective area. In order to resolve the ambiguity,
\ac{PMIx} introduced the concept of a \declareterm{data realm}\emph{data realm}
that identifies the scope to which the referenced data applies. Thus, reference
to an attribute that isn't specific to a realm (e.g., the
\refattr{PMIX_NUM_NODES} attribute) must be accompanied by a corresponding
attribute identifying the realm to which the request pertains if it differs
from the implementation's default.

\ac{PMIx} has defined five \emph{data realms} to resolve the ambiguities, as
captured in the following attributes used in \refapi{PMIx_Get} for retrieving
information from each of the realms:

%
\declareAttribute{PMIX_SESSION_INFO}{"pmix.ssn.info"}{bool}{
Return information regarding the session realm of the target process.
}

%
\declareAttribute{PMIX_JOB_INFO}{"pmix.job.info"}{bool}{
Return information regarding the job realm corresponding to the namespace in
the target process' identifier.
}

%
\declareAttribute{PMIX_APP_INFO}{"pmix.app.info"}{bool}{
Return information regarding the application realm to which the target process
belongs - the namespace of the target process serves to identify the job
containing the target application. If information about an application other
than the one containing the target process is desired, then the attribute array
must contain a \refattr{PMIX_APPNUM} attribute identifying the desired target
application - this is useful when the requesting process knows there are
multiple applications, but doesn't know which processes belong to which
application.
}

\declareAttribute{PMIX_PROC_INFO}{"pmix.proc.info"}{bool}{
Return information regarding the target process. This attribute is technically
not required as the \refapi{PMIx_Get} \ac{API} specifically identifies the
target process in its parameters. However, it is included here for
completeness.
}

%
\declareAttribute{PMIX_NODE_INFO}{"pmix.node.info"}{bool}{
Return information from the node realm regarding the node. If information about
a node other than the one containing the requesting process is desired, then
the attribute array must also contain either the \refattr{PMIX_NODEID} or
\refattr{PMIX_HOSTNAME} attribute identifying the desired target.
}

The \ac{PMIx} server has corresponding attributes the host can use to specify the realm of information that it provides during namespace registration (see Section \ref{chap:api_server:assemble}).

\adviceuserstart
If information about a session other than the one containing the requesting
process is desired, then the attribute array must contain a
\refattr{PMIX_SESSION_ID} attribute identifying the desired target session.
This is required as many environments only guarantee unique namespaces within a
session, and not across sessions.
\adviceuserend

%%%%%%%%%%%%%%%%%%%%%%%%%%%%%%
\subsection{Session realm keys}

Session-related information can only be retrieved by including the \refattr{PMIX_SESSION_INFO} attribute in the \refarg{info} array passed to \refapi{PMIx_Get}. If information about a session other than the one containing the
requesting process is desired, then the \refarg{info} array must contain a
\refattr{PMIX_SESSION_ID} attribute identifying the desired target session.
This is required as many environments only guarantee unique namespaces
within a session, and not across sessions.

Note that the \refarg{proc} argument of \refapi{PMIx_Get} is ignored when
referencing session-related information.

Session-level information includes the following keys:

%
\declareAttribute{PMIX_SESSION_ID}{"pmix.session.id"}{uint32_t}{
Session identifier assigned by the scheduler.
}

%
\declareAttribute{PMIX_CLUSTER_ID}{"pmix.clid"}{char*}{
A string name for the cluster this allocation is on
}

\declareAttribute{PMIX_ALLOCATED_NODELIST}{"pmix.alist"}{char*}{
Comma-delimited list or regular expression of all nodes in this allocation regardless of whether or not they currently host processes.
}

%
\declareAttribute{PMIX_UNIV_SIZE}{"pmix.univ.size"}{uint32_t}{
Number of allocated slots in a session - each slot may or may not be occupied
by an executing process. Note that this attribute is equivalent to providing
the \refattr{PMIX_MAX_PROCS} entry in the \refattr{PMIX_SESSION_INFO_ARRAY}
array - it is included in the Standard for historical reasons.
}

%
\declareAttribute{PMIX_MAX_PROCS}{"pmix.max.size"}{uint32_t}{
Maximum number of processes that can be executed in the specified realm.
Typically, this is a constraint imposed by a scheduler or by user settings in a
hostfile or other resource description.
}

\declareAttribute{PMIX_NODE_LIST}{"pmix.nlist"}{char*}{
Comma-delimited list of nodes assigned to the specified realm.
}
%
\declareAttribute{PMIX_NUM_SLOTS}{"pmix.num.slots"}{uint32_t}{
Number of slots allocated in the specified realm. Note that this attribute is
the equivalent to \refattr{PMIX_MAX_PROCS} used in the corresponding context -
it is included in the Standard for historical reasons.
}

%
\declareAttributeNEW{PMIX_NUM_NODES}{"pmix.num.nodes"}{uint32_t}{
Number of nodes in the specified realm.
}

%
\declareAttribute{PMIX_TMPDIR}{"pmix.tmpdir"}{char*}{
Full path to the top-level temporary directory assigned to the session.
}

%
\declareAttribute{PMIX_TDIR_RMCLEAN}{"pmix.tdir.rmclean"}{bool}{
Resource Manager will cleanup assigned temporary directory trees.
}

Attributes used to describe the \ac{RM} - these are values assigned by the host environment:
%
\declareAttribute{PMIX_RM_NAME}{"pmix.rm.name"}{char*}{
String name of the \ac{RM}.
}

%
\declareAttribute{PMIX_RM_VERSION}{"pmix.rm.version"}{char*}{
\ac{RM} version string.
}

%%%%%%%%%%%%%%%%%%%%%%%%%%%%%%
\subsection{Job realm keys}

Job-related information is retrieved by including the namespace of the target
job and a rank of \refconst{PMIX_RANK_WILDCARD} in the \refarg{proc} argument
passed to \refapi{PMIx_Get}. If desired for code clarity, the caller can also
include the \refattr{PMIX_JOB_INFO} attribute in the \refarg{info} array,
though this is not required. If information is requested about a namespace in
a session other than the one containing the requesting process, then the
\refarg{info} array must contain a \refattr{PMIX_SESSION_ID} attribute
identifying the desired target session. This is required as
many environments only guarantee unique namespaces within a session, and not
across sessions.

Job-level information includes the following keys:

%
\declareAttribute{PMIX_NSPACE}{"pmix.nspace"}{char*}{
Namespace of the job - may be a numerical value expressed as a string, but is often a non-numerical string carrying information solely of use to the system.
}

%
\declareAttribute{PMIX_JOBID}{"pmix.jobid"}{char*}{
Job identifier assigned by the scheduler - may be identical to the namespace, but is often a numerical value expressed as a string (e.g., "12345.3").
}

%
\declareAttribute{PMIX_NPROC_OFFSET}{"pmix.offset"}{pmix_rank_t}{
Starting global rank of this job - referenced using \refconst{PMIX_RANK_WILDCARD}.
}

%
\pasteAttributeItem{PMIX_MAX_PROCS}
%
\pasteAttributeItemBegin{PMIX_NUM_SLOTS}Jobs may reserve a subset of their
assigned slots for dynamic operations such as \refapi{PMIx_Spawn} - i.e., not
all slots may be occupied by executing processes from this job at a given
point in time.
\pasteAttributeItemEnd{}
%
\pasteAttributeItemBegin{PMIX_NUM_NODES}This is the number of nodes assigned
to the job, which may be a subset of the nodes allocated to the overall
session. Jobs may reserve a subset of their assigned nodes for dynamic
operations such as \refapi{PMIx_Spawn} - i.e., not all
nodes may have executing processes from this job at a given point in time.
\pasteAttributeItemEnd{}

\pasteAttributeItem{PMIX_NODE_LIST}

%
\declareAttributeNEW{PMIX_CMD_LINE}{"pmix.cmd.line"}{char*}{
Command line used to execute the specified job (e.g., "mpirun -n 2 --map-by foo ./myapp")
}


%
\declareAttribute{PMIX_NSDIR}{"pmix.nsdir"}{char*}{
Full path to the temporary directory assigned to the namespace, under \refattr{PMIX_TMPDIR}.
}

%
\declareAttribute{PMIX_LOCALLDR}{"pmix.lldr"}{pmix_rank_t}{
Lowest rank on this node within this job - referenced using \refconst{PMIX_RANK_WILDCARD}.
}

%
\declareAttribute{PMIX_JOB_SIZE}{"pmix.job.size"}{uint32_t}{
Total number of processes in this job across all contained applications. Note that this value can be different from \refattr{PMIX_MAX_PROCS}. For example, users may choose to subdivide an allocation (running several jobs in parallel within it), and dynamic programming models may support adding and removing processes from a running \refterm{job} on-the-fly. In the latter case, \ac{PMIx} events must be used to notify processes within the job that the job size has changed.
}

%
\declareAttribute{PMIX_JOB_NUM_APPS}{"pmix.job.napps"}{uint32_t}{
Number of applications in this job.
}


%%%%%%%%%%%%%%%%%%%%%%%%%%%%%%
\subsection{Application realm keys}

Application-related information can only be retrieved by including the
\refattr{PMIX_APP_INFO} attribute in the \refarg{info} array passed to
\refapi{PMIx_Get}.

Application-level information includes the following keys:

\pasteAttributeItem{PMIX_APPNUM}

%
\pasteAttributeItem{PMIX_MAX_PROCS}
%
\pasteAttributeItem{PMIX_NUM_SLOTS}
%
\pasteAttributeItemBegin{PMIX_NUM_NODES}This is the number of nodes allocated
to the job, which may be a subset of the nodes allocated to the overall
session. Jobs may utilize a subset of their allocated nodes - i.e., not all
nodes may have executing processes from this job at a given point in time.
\pasteAttributeItemEnd{}

%
\declareAttribute{PMIX_APPLDR}{"pmix.aldr"}{pmix_rank_t}{
Lowest rank in this application within this job - referenced using \refconst{PMIX_RANK_WILDCARD}.
}

%
\declareAttribute{PMIX_APP_SIZE}{"pmix.app.size"}{uint32_t}{
Number of processes in this application.
}

%
\declareAttributeNEW{PMIX_APP_ARGV}{"pmix.app.argv"}{char*}{
Consolidated argv passed to the spawn command for the given application (e.g., "./myapp arg1 arg2 arg3")
}

%%%%%%%%%%%%%%%%%%%%%%%%%%%%%%
\subsection{Process realm keys}

Process-related information is retrieved by simply referencing the namespace
and rank of the target process in the call to \refapi{PMIx_Get}. If information
is requested about a process in a session other than the one containing the
requesting process, then an attribute identifying the target session must be
provided. This is required as many environments only guarantee unique
namespaces within a session, and not across sessions.

Process-level information includes the following keys:

%
\declareAttribute{PMIX_APPNUM}{"pmix.appnum"}{uint32_t}{
The application number within the job in which this process is a member.
}

%
\declareAttribute{PMIX_RANK}{"pmix.rank"}{pmix_rank_t}{
Process rank within the job.
}

%
\declareAttribute{PMIX_GLOBAL_RANK}{"pmix.grank"}{pmix_rank_t}{
Process rank spanning across all jobs in this session.
}

%
\declareAttribute{PMIX_APP_RANK}{"pmix.apprank"}{pmix_rank_t}{
Process rank within this application.
}

%
\declareAttribute{PMIX_PROC_URI}{"pmix.puri"}{char*}{
\ac{URI} containing contact information for a given process.
}

%
\declareAttribute{PMIX_PARENT_ID}{"pmix.parent"}{pmix_proc_t}{
Process identifier of the parent process of the calling process.
}

%
\declareAttribute{PMIX_EXIT_CODE}{"pmix.exit.code"}{int}{
Exit code returned when process terminated
}

%
\declareAttribute{PMIX_PROCID}{"pmix.procid"}{pmix_proc_t}{
Process identifier
}

%
\declareAttribute{PMIX_LOCAL_RANK}{"pmix.lrank"}{uint16_t}{
Local rank on this node within this job.
}

%
\declareAttribute{PMIX_NODE_RANK}{"pmix.nrank"}{uint16_t}{
Process rank on this node spanning all jobs.
}

%
\declareAttributeNEW{PMIX_PACKAGE_RANK}{"pmix.pkgrank"}{uint16_t}{
Process rank within this job on the package where this process resides
}

%
\declareAttribute{PMIX_PROC_PID}{"pmix.ppid"}{pid_t}{
\ac{PID} of specified process.
}

%
\declareAttribute{PMIX_PROCDIR}{"pmix.pdir"}{char*}{
Full path to the subdirectory under \refattr{PMIX_NSDIR} assigned to the process.
}

%
\declareAttribute{PMIX_CPUSET}{"pmix.cpuset"}{char*}{
HWLOC\footnote{\url{https://www.open-mpi.org/projects/hwloc/}} bitmap applied to the process upon launch.
}

%
\declareAttribute{PMIX_CREDENTIAL}{"pmix.cred"}{char*}{
Security credential assigned to the process.
}

%
\declareAttribute{PMIX_SPAWNED}{"pmix.spawned"}{bool}{
\code{true} if this process resulted from a call to \refapi{PMIx_Spawn}. Lack of inclusion (i.e., a return status of \refconst{PMIX_ERR_NOT_FOUND}) corresponds to a value of \code{false} for this attribute.
}

%
\declareAttribute{PMIX_ARCH}{"pmix.arch"}{uint32_t}{
Architecture flag.
}

%
\declareAttributeNEW{PMIX_REINCARNATION}{"pmix.reinc"}{uint32_t}{
Number of times this process has been re-instantiated - i.e, a value of zero indicates that the process has never failed and been restarted.
}

%
\declareAttribute{PMIX_CRED_TYPE}{"pmix.sec.ctype"}{char*}{
When passed in \refapi{PMIx_Get_credential}, a prioritized, comma-delimited list of desired credential types for use
in environments where multiple authentication mechanisms may be available. When returned in a callback function, a
string identifier of the credential type.
}

%
\declareAttribute{PMIX_CRYPTO_KEY}{"pmix.sec.key"}{pmix_byte_object_t}{
Blob containing crypto key
}

%%%%%%%%%%%%%%%%%%%%%%%%%%%%%%
\subsection{Node realm keys}


Node-level information includes the following keys:

%
\declareAttribute{PMIX_HOSTNAME}{"pmix.hname"}{char*}{
Name of the host (e.g., where a specified process is running, or a given device is located).
}

%
\declareAttributeNEW{PMIX_HOSTNAME_ALIASES}{"pmix.alias"}{char*}{
Comma-delimited list of names by which this node is known.
}

%
\declareAttributeNEW{PMIX_HOSTNAME_KEEP_FQDN}{"pmix.fqdn"}{bool}{
FQDN hostnames are being retained
}

%
\declareAttribute{PMIX_NODEID}{"pmix.nodeid"}{uint32_t}{
Node identifier expressed as the node's index (beginning at zero) in an array of nodes within the active session. The value must be unique and directly correlate to the \refattr{PMIX_HOSTNAME} of the node - i.e., users can interchangeably reference the same location using either the \refattr{PMIX_HOSTNAME} or corresponding \refattr{PMIX_NODEID}.
}

%
\declareAttribute{PMIX_LOCAL_PEERS}{"pmix.lpeers"}{char*}{
Comma-delimited list of ranks on this node within the specified namespace - referenced using \refconst{PMIX_RANK_WILDCARD}.
}

%
\declareAttribute{PMIX_LOCAL_PROCS}{"pmix.lprocs"}{pmix_proc_t array}{
Array of \refstruct{pmix_proc_t} of all processes on the specified node - referenced using \refconst{PMIX_RANK_WILDCARD}.
}

%
\declareAttribute{PMIX_LOCAL_CPUSETS}{"pmix.lcpus"}{char*}{
Colon-delimited cpusets of local peers within the specified namespace - referenced using \refconst{PMIX_RANK_WILDCARD}.
}

%
\declareAttribute{PMIX_LOCAL_SIZE}{"pmix.local.size"}{uint32_t}{
Number of processes in this job or application on this node.
}

%
\declareAttribute{PMIX_NODE_SIZE}{"pmix.node.size"}{uint32_t}{
Number of processes across all jobs on this node.
}

\section{Retrieval rules for reserved keys}
The retrieval rules for reserved keys are relatively simple as the keys are
required, by definition, to be available when the client begins execution.
Accordingly, \refapi{PMIx_Get} for a reserved key first checks the local
\ac{PMIx} Client cache for the target key. If the information is not found,
then the \refconst{PMIX_ERR_NOT_FOUND} error constant is returned unless the
target process belongs to a different namespace from that of the requester.

In the case where the target and requester's namespaces differ, then the
request is forwarded to the local \ac{PMIx} server. Upon receiving the
request, the server shall check its data storage for the specified namespace.
If it already knows about this namespace, then it shall attempt to lookup the
specified key, returning the value if it is found or the
\refconst{PMIX_ERR_NOT_FOUND} error constant.

If the server does not have a copy of the information for the specified
namespace, then the server shall take one of the following actions:
\begin{enumerate}
    \item If the request included the \refattr{PMIX_IMMEDIATE} attribute, then
    the server will respond to the client with the
    \refconst{PMIX_ERR_NOT_FOUND} status

    \item If the host has provided the \ac{DBCX} module function interface
    (\refapi{pmix_server_dmodex_req_fn_t}), then the server shall pass the
    request to its host for servicing. The host is responsible for identifying
    a source of information on the specified namespace and retrieving it. The
    host is required to retrieve \emph{all} of the namespace-level information
    and returning it to the requesting server in anticipation of further
    requests regarding the target namespace. If the host cannot retrieve the
    namespace information, then it must respond with the \refconst{PMIX_ERR_NOT_FOUND} error constant unless the \refattr{PMIX_TIMEOUT} is given and reached (in which case, the host must respond with the \refconst{PMIX_ERR_TIMEOUT} constant).

    \item If the host does not support the \ac{DBCX} interface, then the
    server will respond to the client with the \refconst{PMIX_ERR_NOT_FOUND}
    status
\end{enumerate}

The \refterm{data realm} that is to be searched for the requested information is communicated through a combination of the target process ID (i.e., the \refarg{proc} argument to refapi{PMIx_Get}) and attributes passed via the \refarg{info} array, according to the following rules:

\begin{itemize}
    \item Session-level information shall be referenced using a
    \item Job-related information shall be referenced using a \refarg{proc} that specifies the namespace of the target job with a rank of \refconst{PMIX_RANK_WILDCARD}.
\end{itemize}

RHC: SPECIAL HANDLING REQUIRED TO SUPPORT REALM-SPECIFIC INFO - NEED TO EXPLAIN WHICH KEYS ARE MAPPED TO WHICH REALMS, HOW THE CLIENT LIBRARY DEALS WITH THEM, DETECTION OF REALM DIRECTIVES (E.G., PMIX_APP_INFO) AND HOW THAT IMPACTS THE SEARCH

%%%%%%%%%%%
\subsection{Accessing information: examples}
\label{chap:api_kv:getex}

This section provides examples illustrating methods for accessing information at various levels. The intent of the examples is not to provide comprehensive coding guidance, but rather to illustrate how \refapi{PMIx_Get} can be used to obtain information on a \refterm{session}, \refterm{job}, \refterm{application}, process, and node.

\subsubsection{Session-level information}

The \refapi{PMIx_Get} \ac{API} does not include an argument for specifying the \refterm{session} associated with the information being requested. Information regarding the session containing the requestor can be obtained by the following methods:

\begin{itemize}
\item for session-level attributes (e.g., \refattr{PMIX_UNIV_SIZE}), specifying the requestor's namespace and a rank of \refconst{PMIX_RANK_WILDCARD}; or
\item for non-specific attributes (e.g., \refattr{PMIX_NUM_NODES}), including the \refattr{PMIX_SESSION_INFO} attribute to indicate that the session-level information for that attribute is being requested
\end{itemize}

Example requests are shown below:

\cspecificstart
\begin{codepar}
pmix_info_t info;
pmix_value_t *value;
pmix_status_t rc;
pmix_proc_t myproc, wildcard;

/* initialize the client library */
PMIx_Init(&myproc, NULL, 0);

/* get the #slots in our session */
PMIX_PROC_LOAD(&wildcard, myproc.nspace, PMIX_RANK_WILDCARD);
rc = PMIx_Get(&wildcard, PMIX_UNIV_SIZE, NULL, 0, &value);

/* get the #nodes in our session */
PMIX_INFO_LOAD(&info, PMIX_SESSION_INFO, NULL, PMIX_BOOL);
rc = PMIx_Get(&wildcard, PMIX_NUM_NODES, &info, 1, &value);
\end{codepar}
\cspecificend

Information regarding a different session can be requested by either specifying the namespace and a rank of \refconst{PMIX_RANK_WILDCARD} for a process in the target session, or adding the \refattr{PMIX_SESSION_ID} attribute identifying the target session. In the latter case, the \refarg{proc} argument to \refapi{PMIx_Get} will be ignored:

\cspecificstart
\begin{codepar}
pmix_info_t info[2];
pmix_value_t *value;
pmix_status_t rc;
pmix_proc_t myproc;
uint32_t sid;

/* initialize the client library */
PMIx_Init(&myproc, NULL, 0);

/* get the #nodes in a different session */
sid = 12345;
PMIX_INFO_LOAD(&info[0], PMIX_SESSION_INFO, NULL, PMIX_BOOL);
PMIX_INFO_LOAD(&info[1], PMIX_SESSION_ID, &sid, PMIX_UINT32);
rc = PMIx_Get(&myproc, PMIX_NUM_NODES, info, 2, &value);
\end{codepar}
\cspecificend

\subsubsection{Job-level information}

Information regarding a job can be obtained by the following methods:

\begin{itemize}
\item for job-level attributes (e.g., \refattr{PMIX_JOB_SIZE} or \refattr{PMIX_JOB_NUM_APPS}), specifying the namespace of the job and a rank of \refconst{PMIX_RANK_WILDCARD} for the \refarg{proc} argument to \refapi{PMIx_Get}; or
\item for non-specific attributes (e.g., \refattr{PMIX_NUM_NODES}), including the \refattr{PMIX_JOB_INFO} attribute to indicate that the job-level information for that attribute is being requested
\end{itemize}

Example requests are shown below:

\cspecificstart
\begin{codepar}
pmix_info_t info;
pmix_value_t *value;
pmix_status_t rc;
pmix_proc_t myproc, wildcard;

/* initialize the client library */
PMIx_Init(&myproc, NULL, 0);

/* get the #apps in our job */
PMIX_PROC_LOAD(&wildcard, myproc.nspace, PMIX_RANK_WILDCARD);
rc = PMIx_Get(&wildcard, PMIX_JOB_NUM_APPS, NULL, 0, &value);

/* get the #nodes in our job */
PMIX_INFO_LOAD(&info, PMIX_JOB_INFO, NULL, PMIX_BOOL);
rc = PMIx_Get(&wildcard, PMIX_NUM_NODES, &info, 1, &value);
\end{codepar}
\cspecificend


\subsubsection{Application-level information}

Information regarding an application can be obtained by the following methods:

\begin{itemize}
\item for application-level attributes (e.g., \refattr{PMIX_APP_SIZE}), specifying the namespace and rank of a process within that application;
\item for application-level attributes (e.g., \refattr{PMIX_APP_SIZE}), including the \refattr{PMIX_APPNUM} attribute specifying the application whose information is being requested. In this case, the namespace field of the \refarg{proc} argument is used to reference the \refterm{job} containing the application - the \refterm{rank} field is ignored;
\item or application-level attributes (e.g., \refattr{PMIX_APP_SIZE}), including the \refattr{PMIX_APPNUM} and \refattr{PMIX_NSPACE} or \refattr{PMIX_JOBID} attributes specifying the job/application whose information is being requested. In this case, the \refarg{proc} argument is ignored;
\item for non-specific attributes (e.g., \refattr{PMIX_NUM_NODES}), including the \refattr{PMIX_APP_INFO} attribute to indicate that the application-level information for that attribute is being requested
\end{itemize}

Example requests are shown below:

\cspecificstart
\begin{codepar}
pmix_info_t info;
pmix_value_t *value;
pmix_status_t rc;
pmix_proc_t myproc, otherproc;
uint32_t appsize, appnum;

/* initialize the client library */
PMIx_Init(&myproc, NULL, 0);

/* get the #processes in our application */
rc = PMIx_Get(&myproc, PMIX_APP_SIZE, NULL, 0, &value);
appsize = value->data.uint32;

/* get the #nodes in an application containing "otherproc".
 * Note that the rank of a process in the other application
 * must be obtained first - a simple method is shown here */

/* assume for this example that we are in the first application
 * and we want the #nodes in the second application - use the
 * rank of the first process in that application, remembering
 * that ranks start at zero */
PMIX_PROC_LOAD(&otherproc, myproc.nspace, appsize);

PMIX_INFO_LOAD(&info, PMIX_APP_INFO, NULL, PMIX_BOOL);
rc = PMIx_Get(&otherproc, PMIX_NUM_NODES, &info, 1, &value);

/* alternatively, we can directly ask for the #nodes in
 * the second application in our job, again remembering that
 * application numbers start with zero */
appnum = 1;
PMIX_INFO_LOAD(&appinfo[0], PMIX_APP_INFO, NULL, PMIX_BOOL);
PMIX_INFO_LOAD(&appinfo[1], PMIX_APPNUM, &appnum, PMIX_UINT32);
rc = PMIx_Get(&myproc, PMIX_NUM_NODES, appinfo, 2, &value);

\end{codepar}
\cspecificend

\subsubsection{Process-level information}

Process-level information is accessed by providing the namespace and rank of the target process. In the absence of any directive as to the level of information being requested, the \ac{PMIx} library will always return the process-level value.

\subsubsection{Node-level information}

Information regarding a node within the system can be obtained by the following methods:

\begin{itemize}
\item for node-level attributes (e.g., \refattr{PMIX_NODE_SIZE}), specifying the namespace and rank of a process executing on the target node;
\item for node-level attributes (e.g., \refattr{PMIX_NODE_SIZE}), including the \refattr{PMIX_NODEID} or \refattr{PMIX_HOSTNAME} attribute specifying the node whose information is being requested. In this case, the \refarg{proc} argument's values are ignored; or
\item for non-specific attributes (e.g., \refattr{PMIX_MAX_PROCS}), including the \refattr{PMIX_NODE_INFO} attribute to indicate that the node-level information for that attribute is being requested
\end{itemize}

Example requests are shown below:

\cspecificstart
\begin{codepar}
pmix_info_t info[2];
pmix_value_t *value;
pmix_status_t rc;
pmix_proc_t myproc, otherproc;
uint32_t nodeid;

/* initialize the client library */
PMIx_Init(&myproc, NULL, 0);

/* get the #procs on our node */
rc = PMIx_Get(&myproc, PMIX_NODE_SIZE, NULL, 0, &value);

/* get the #slots on another node */
PMIX_INFO_LOAD(&info[0], PMIX_NODE_INFO, NULL, PMIX_BOOL);
PMIX_INFO_LOAD(&info[1], PMIX_HOSTNAME, "remotehost", PMIX_STRING);
rc = PMIx_Get(&myproc, PMIX_MAX_PROCS, info, 2, &value);

\end{codepar}
\cspecificend

